% Nouveau EGA V (cote 25): archive cote 25, site page 119 (PDF page 120).
% typescript. Draft compared: none.
% AI-transcribed by Claude Opus 5.5 (claude-opus-5-5, adaptive thinking, effort high), Message Batches API, batch ega25b (msgbatch_01A7n5YuXubUPNK1zsepi8RG), 2026-11-13; checked against both supplied scan parts; see CHECK_DECISIONS.tsv.
\scanpage{25/119}
\var{\margin{\unsure{22|2/2}}}{\placedmargin{\unsure{22|2/2}}}{marginal note placed}{ega25_25_119-v0001}
\var{\margin{$|$}}{\placedmargin{$|$}}{marginal note placed}{ega25_25_119-v0002}
Je me répens d'avoir donné la proposition \var{\deleted{\unsure{3|2}}}{}{deletion}{ega25_25_119-v0003} sous une forme éculée et
pisseuse, à mi chemin entre  les hypothèses classiques et les hypothèses
naturelles, et sans donner une récip\unsure{r|p}oque. Aussi je propose l'énoncé:
\uline{Proposition\hand{3}\unsure{\var{\over{2}{3}}{3}{writing over}{ega25_25_119-v0004}|3}. \var{\deleted{revue et corrigée}}{}{deletion}{ega25_25_119-v0005}}.  Soit $X \hand{\to} S$ plat loc de présentation
finie,  $Q=\overset{\hand{\vee}}{P} = P(\overset{\hand{\vee}}{E})$  un fibré \var{\xout{vectoriel sur}}{}{typewriter strike-out}{ega25_25_119-v0006} projectif sur $S$ défini par
un Module loc  libre de type fini $F = \overset{\hand{\vee}}{E}$ ,  considérons l'ensemble
$\hand{\unsure{\phi|\emptyset}}$  des systêmes linéaires \var{\inserted[typed]{$D$ \var{\deleted{\unsure{E}}}{}{deletion}{ega25_25_119-v0008}}}{$D$ \var{\deleted{\unsure{E}}}{}{deletion}{ega25_25_119-v0008}}{insertion}{ega25_25_119-v0007} de diviseurs \var{\xout{sur}}{}{typewriter strike-out}{ega25_25_119-v0009} sur $X$, paramétrés par $Q$, tels
que l'ensemble  $Z$ des points fixes de $D$  satisfasse la propriété: $z \hand{\in} Z$
implique prof$_z \underline{O}_{X_s} \hand{\geqslant} 2$  (où  $s$ \var{\xout{\unsure{$= g(z)$)}}}{}{typewriter strike-out}{ega25_25_119-v0010} est l'image de $z$ dans $S$). \var{\xout{Alors}}{}{typewriter strike-out}{ega25_25_119-v0011}
\var{\xout{\unsure{$D \hand{\in} \hand{\unsure{\phi|\emptyset}}$ associons le morphisme correspondant $f$ $X$-$Z$ $P$ considérée}}}{}{typewriter strike-out}{ega25_25_119-v0012}
\var{\xout{\unsure{considérée comme}}}{}{typewriter strike-out}{ega25_25_119-v0013}  Soit $\hand{\unsure{\psi|\emptyset}}$ l'ensemble des  \var{\deleted{``applications rationelles}}{}{deletion}{ega25_25_119-v0014}\var{\inserted{pseudo-morphismes}}{pseudo-morphismes}{insertion}{ega25_25_119-v0015} rel à
$S$ \var{\deleted{\unsure{*}}}{}{deletion}{ega25_25_119-v0016}\var{\inserted[typed]{$f$}}{$f$}{insertion}{ega25_25_119-v0017} de $X$ dans  $P$, telles que le\var{\xout{\illegible{1}}}{}{typewriter strike-out}{ega25_25_119-v0018} domaine de définition \var{\inserted{$U=$}}{$U=$}{insertion}{ega25_25_119-v0019}\var{\over{dom}{\unsure{U}}}{\unsure{U}}{writing over}{ega25_25_119-v0020}(f)
\var{\xout{\unsure{f satisf}}}{}{typewriter strike-out}{ega25_25_119-v0021} satisfasse la condition : $z \hand{\in} X$-$U$\var{\deleted{(f)}}{}{deletion}{ega25_25_119-v0022} $\hand{\Longrightarrow}$ prof$_z \underline{O}_{X_s} \hand{\geqslant} 2$\var{\inserted{$\unsure{\surd}$, et $f_U = f|U : U \to \check{P}$ satisfait la condition de non dégénérescence envisagée dans 3.1.}}{$\unsure{\surd}$, et $f_U = f|U : U \to \check{P}$ satisfait la condition de non dégénérescence envisagée dans 3.1.}{insertion}{ega25_25_119-v0023}. Considé\var{\xout{rons}}{}{typewriter strike-out}{ega25_25_119-v0024}
rons l'application naturelle  $D \hand{\rightsquigarrow} f_D$  de $\hand{\unsure{\phi|\emptyset}}$  dans $\hand{\unsure{\psi|\emptyset}}$ .  Alors :

a)  Cette application est injective,  et  pour $D \hand{\in} \hand{\unsure{\phi|\emptyset}}$ , l'ensemble des
points fixes $Z$  n'est autre que l'ensemble  complémentaire du domaine de
définition \var{\deleted{\unsure{Dif}}}{}{deletion}{ega25_25_119-v0025} $U$\var{\deleted{(f)}}{}{deletion}{ega25_25_119-v0026} de $f$\hand{$_D$}.

b)  \var{\xout{\unsure{Pour qu'une f $\hand{\unsure{\psi|\emptyset}}$ appartienne à provienne d'une $D$ $\hand{\unsure{\phi|\emptyset}}$, il}}}{}{typewriter strike-out}{ega25_25_119-v0027} Soit
$f \hand{\in} \hand{\unsure{\psi|\emptyset}}$ , et soit $U$ un ouvert sur lequel $f$ soit défini, et tel que $z \hand{\in} X$-$U$
implique prof $\underline{O}_{X_s,z} \hand{\geqslant} \unsure{2|0}$ \var{\xout{(il en existe}}{}{typewriter strike-out}{ega25_25_119-v0028} (par exemple $U = U(f)$ domaine de
définition de $f$). Pour que $f$ provienne d'une $D \hand{\in} \hand{\unsure{\phi|\emptyset}}$ , il f et suffit  que\hand{\unsure{$\surd$}}
\var{\deleted{$f_U$ satisfasse la condition énoncée dans 3.1. et que}}{}{deletion}{ega25_25_119-v0029}
posant $\underline{L}_U = f_U^{\hand{\unsure{*|\bullet}}}(\underline{O}_P(1))$  (où  $f_U$: $U \hand{\to} P$  est le morphisme induit par $f$),
\var{\deleted{$f_U$ satisfasse la condition énoncée dans 3.1 et}}{}{deletion}{ega25_25_119-v0030}
\var{\xout{en ait}}{}{typewriter strike-out}{ega25_25_119-v0031} le Module  $i_{\hand{*}}(\underline{L}_U)$  sur $X$ soit inversible (où $i$:$U \hand{\to} X$ est l'immer-%
sion canonique).

On remarquera que lorsque  les fibres de $X$ sur $S$ satisfont à $(S_2)$,
par exemple sont normales, voire géométriquement  normales, alors la
condition de profondeur  envisagée sur  une partie fermée $Z$ de $X$ dans \var{\deleted{\illegible{1}}}{}{deletion}{ega25_25_119-v0032}
la proposition  signifie simplement que pour tout $s \hand{\in} S$,

\mspage{119}
